\label{sec:appendix-configuration}
\label{sec:configuration-reference}
\label{app:configuration-reference}

This appendix lists the configuration settings registered by BioseqDB.
There are 73 production settings in the main extension and two in the
optional \texttt{bioseq\_cache} extension. Three further controls are compiled
only into test builds and appear in a separate table. The reference also
identifies every field of \texttt{aligner\_options}, including the index-read
requirement that has no configuration-setting equivalent.

\section{Reading the reference}

A PostgreSQL configuration setting is also called a \emph{GUC}. Each table
uses its full name, so a row can be used directly with \texttt{SHOW} or
\texttt{SET}. The built-in default is \textbf{bold} and appears first.
For an enum or Boolean setting, each value has its own description.
The default row also explains the setting's overall purpose.
Compression depends on the compiled codecs. Its table puts Zstandard
first because it is the preferred default when available. The LZ4 and
PGLZ rows explain which default applies when it is unavailable.
An integer or real row gives the inclusive range accepted by \texttt{SET}.
The type abbreviations are \texttt{int} for integer, \texttt{real} for a
fractional number, \texttt{bool} for Boolean and \texttt{enum} for a named choice.

The three alignment algorithm GUCs use names in \texttt{SET}, but the
corresponding \texttt{aligner\_options} fields use integer codes. Their value
rows give both forms. Other GUC enum ordinals are implementation details and
are not per-search option codes.

The defaults describe a newly registered setting. Server configuration,
database settings, role settings and earlier session commands can replace
them. \texttt{RESET} restores the session's reset value, which need not be
the built-in default. \texttt{SET LOCAL} changes a setting only until the
current transaction ends. PostgreSQL documents these distinctions in its
configuration and \texttt{SET} references
\autocite{postgresql17_setting_parameters,postgresql17_set}.

\begin{lstlisting}[style=bioseq-search-sql,
  caption={Inspect a registered setting and compare its current and built-in values.},
  label={lst:configuration-inspect}]
SHOW bioseq.aligner.minimizer.algorithm;

SELECT name, setting, boot_val, reset_val,
       vartype, context, min_val, max_val, enumvals
FROM pg_settings
WHERE name = 'bioseq.aligner.minimizer.algorithm';
\end{lstlisting}

The extension library must have been loaded in the current backend before
all its settings are visible. Settings marked \emph{hidden} are omitted
from \texttt{SHOW ALL} and \texttt{pg\_settings}, but their exact names can
still be read with \texttt{SHOW}. The source inventory therefore includes
more rows than a query of \texttt{pg\_settings} alone
\autocite{postgresql17_pg_settings,postgresql17_guc_functions}.

Unmarked settings can be changed by an ordinary user. A \emph{superuser}
setting requires superuser authority or a granted PostgreSQL permission to
set that parameter. A \emph{reload} setting changes through server
configuration and a configuration reload, rather than ordinary session
\texttt{SET}. MiB means $2^{20}$ bytes and corresponds to PostgreSQL's
\texttt{MB} unit. GiB means $2^{30}$ bytes. The shard limit takes a plain
integer GiB count, not a GUC memory-unit suffix.

An accepted setting is not a promise that every combination can execute.
Algorithm support depends on the alphabet, compiled profiles, hardware and
requested output. For example, disabling the custom scan does not turn
\texttt{mmsearch} into an ordinary scalar search function. The tables identify
controls whose values are stored but currently have no execution consumer.

\section{Relationships between settings}

Maintenance, reading and cache reuse are distinct operations.
Section~\ref{sec:index-update-policy} explains the stored per-index
maintenance policies. Section~\ref{sec:index-read-policy} explains why the
per-search requirement takes precedence over optional index reading.
The default grouped builder publishes an immutable generation explicitly.
It cannot use ordinary insertion maintenance or the older asynchronous
writer to maintain an MPF reader file.

The asynchronous switches do not guarantee that a query leaves a saved
index. The current code submits only ordinary MPS indexes. It requires
$k=15,w=10$ and excludes default bounded ntHash. It needs an eligible
on-insert policy and a worker queue. Its selection prefers MPF before MPS,
but the submission path cannot persist MPF. Grouped explicit builds require
both switches off. Section~\ref{sec:index-build-settings} gives the full
distinction.

Table options and process resources also have different lifetimes.
Compression and stripe defaults are copied into newly created table options.
Changing them does not rewrite an existing table. The decompressed column
cache is local to a backend and is normally cleared at scan teardown.
Index caches and the optional query-result cache hold different objects.
A cache's capacity is not a limit on all the memory used by a query.

For the optional result cache, a non-disabled mode requires
\texttt{bioseq\_cache} in \nolinkurl{shared_preload_libraries} and a server
restart. Completed-result age is measured from completion. Reuse applies
the smaller of the producer's time-to-live and the current reader's limit.
The Bioseq alignment adapter currently supports exact request reuse.
Selecting \texttt{reusable\_fragments} does not add partial-range reuse to
that adapter.
